% Čest: teorija-modelov
% Glava: modeli-aritmetiky
% Odděl: izčislimy-modeli

\documentclass[../../../include/open-logic-section]{subfiles}

\begin{document}

\olfileid[isv]{mod}{mar}{cmp}
\section{Izčislimy modeli aritmetiky}

\begin{explain}
Standardny model~$\Struct{N}$ imaje dva prijetna svojstva. Jego domenom
sųt prirodne čisla~$\Nat$, t.j. jego elementi sųt upravo takove
objekty, o ktoryh hočemo govoriti jezykom aritmetiky, a standardny
numeral~$\num{n}$ v samoj stvari označaje~$n$. Drugo prijetno svojstvo
jest, že vse interpretacije nelogičnyh simbolov jezyka~$\Lang{L_A}$
sųt \emph{izčislime}. Funkcije naslědnika, dodavanja i množenja, ktore
služet kako $\Assign{\prime}{N}$, $\Assign{+}{N}$ i
$\Assign{\times}{N}$, sųt izčislime funkcije na čislah. (Jih možno
daže opreděliti s pomočju primitivnoj rekursije, a zato izčisljati
Turingovymi mašinami.) A odnošenje ``menše od'' v~$\Struct{N}$, t.j.
$\Assign{<}{N}$, jest rěšimo.

Nestandardne modeli aritmetičnyh teorij, takovyh kako $\Th{Q}$ i
$\Th{PA}$, morajut sodržati nestandardne elementy. Zato jih često
prědstavljajut s domenami, ktore sodržet kopiju~$\Nat$ i takože
nestandardne !!{element}y. Však vsaka !!{structure} s izčisljivo
bezkonečnym domenom imaje izomorfnu kopiju s libo-ktorym
!!{denumerable} množstvom kako domenom, vključno s~$\Nat$.
Sučstvujut zato takože nestandardne modeli s domenom~$\Nat$. V takovyh
modelah~$\Struct{M}$ něktore čisla, razuměje se, ne mogut igrati svoje
obyčne role, ibo mora sučstvovati něktoro~$k$ različno
od~$\Value{\num{n}}{M}$ za vse~$n \in \Nat$.
\end{explain}

\begin{defn}
!!^a{structure}~$\Struct{M}$ za jezyk~$\Lang{L_A}$ jest
\emph{izčislima} togda i samo togda, kogda $\Domain{M} = \Nat$,
$\Assign{\prime}{M}$, $\Assign{+}{M}$ i $\Assign{\times}{M}$ sųt
izčislime funkcije, a $\Assign{<}{M}$ jest rěšimo odnošenje.
\end{defn}

\begin{ex}\ollabel{ex:comp-model-q}
Pripomnite si strukturu~$\Struct{K}$ iz
\olref[mdq]{ex:model-K-of-Q}. Jej domenom byl $\Domain{K} = \Nat
\cup \{a\}$, a interpretacije byly
\begin{align*}
  \Assign{\Obj{0}}{K} & = 0\\
  \Assign{\prime}{K}(x) & =
  \begin{cases}
    x+1 & \text{ako $x\in \Nat$}\\
    a & \text{ako $x = a$}
  \end{cases}\\
  \Assign{+}{K}(x, y) & =
  \begin{cases}
    x+y & \text{ako $x$, $y \in\Nat$}\\
    a & \text{inače}
  \end{cases}\\
  \Assign{\times}{K}(x, y) & =
  \begin{cases}
    xy & \text{ako $x$, $y \in\Nat$}\\
    0 & \text{ako $x=0$ ili $y=0$}\\
    a & \text{inače}\\
  \end{cases}\\
  \Assign{<}{K} & =
  \Setabs{\tuple{x,y}}{x, y \in \Nat \text{ i } x<y} \cup
  \Setabs{\tuple{x,a}}{x \in \Domain{K}}
\end{align*}
No $\Domain{K}$ jest !!{denumerable}, a zato imaje toliko že elementov,
koliko~$\Nat$. Na priměr, $g\colon \Nat \to \Domain{K}$, pri čem $g(0) =
a$ i $g(n) = n-1$ za $n>0$, jest !!a{bijection}. Možemo prěnesti
interpretacije po~g i tym učiniti~g izomorfizmom medžu novym
modelom~$\Struct{K'}$ teorije~$\Th{Q}$ i~$\Struct{K}$. V~$\Struct{K'}$
moramo simbolam jezyka~$\Lang{L_A}$ pripisati druge funkcije i
odnošenja, ibo druge !!{element}y množstva~$\Nat$ igrajut role
standardnyh i nestandardnyh čisl.

Točno govoręči, $0$ teper igra rolu~$a$, a ne najmenšego standardnogo
čisla. Najmenše standardno čislo teper jest~$1$. Zato postavljajemo
$\Assign{\Obj{0}}{K'} = 1$. Funkcija naslědnika teper takože jest
drugačša: dla standardnogo čisla, t.j. dla $n > 0$, ona po-prěžnjemu
dava $n+1$. No $0$ teper igra rolu~$a$, ktory jest svojim vlastnym
naslědnikom. Zato $\Assign{\prime}{K'}(0) = 0$. Dla dodavanja i
množenja podobno dostavajemo
\begin{align*}
\Assign{+}{K'}(x, y) & =
  \begin{cases}
    x+y-1 & \text{ako $x$, $y >0$}\\
    0 & \text{inače}
  \end{cases}\\
  \Assign{\times}{K'}(x, y) & =
  \begin{cases}
    1 & \text{ako $x = 1$ ili $y = 1$}\\
    xy - x - y + 2 & \text{ako $x$, $y > 1$}\\
    0 & \text{inače}\\
  \end{cases}
\end{align*}
Nadalje, $\tuple{x, y} \in \Assign{<}{K'}$ togda i samo togda, kogda
$x < y$, $x > 0$ i $y > 0$, ili kogda $y = 0$.

Vse te funkcije sųt izčislime funkcije prirodnyh čisl, a
$\Assign{<}{K'}$ jest rěšimo odnošenje na~$\Nat$---no one ne sųt te
same funkcije kako naslědnik, dodavanje i množenje na~$\Nat$, a
$\Assign{<}{K'}$ ne jest tym samym odnošenjem kako~$<$ na~$\Nat$.
\end{ex}

\begin{prob}
Zadajte !!a{structure}~$\Struct{L'}$ s $\Domain{L'} = \Nat$, izomorfnu
strukturi~$\Struct{L}$ iz \olref[mod][mar][mdq]{ex:model-L-of-Q}.
\end{prob}

\begin{explain}
\Olref{ex:comp-model-q} pokazuje, že $\Th{Q}$ imaje izčislime
nestandardne modele s domenom~$\Nat$. Naslědujuča teorema však
pokazuje, že to ne važi za modele teorije~$\Th{PA}$ (a zato takože ne
za modele teorije~$\Th{TA}$).
\end{explain}

\begin{thm}[Tennenbaumova teorema]
$\Struct{N}$ jest, do izomorfizma, jediny izčislimy model
teorije~$\Th{PA}$.
\end{thm}

\end{document}
